\section{引言}

\writingguide{建议写 600--800 字。说明现代 Web 信息规模、用户需求和大语言模型带来的变化；提出本文关注的问题；交代文章组织结构。不要在这里写成摘要。}

Web 已从以超链接和静态页面为主的信息空间，发展为连接用户、服务、知识与智能体的复杂基础设施。搜索系统不仅需要返回与查询相关的页面，还要理解用户意图、组织多模态证据，并支持生成式模型给出可追溯的回答。检索增强生成（Retrieval-Augmented Generation, RAG）将外部信息检索与生成模型结合，为缓解知识过时和事实幻觉提供了一条重要路径\parencite{lewis2020rag}。

\placeholder{继续说明选题背景、研究意义以及本文要回答的核心问题。}

本文围绕 WWW 2026 的“\TrackNameCN”研究方向展开。首先梳理该 Track 的研究边界与技术演进；随后归纳会议论文反映出的主要研究主题；在此基础上，对一篇代表论文进行深入分析；最后讨论当前方法的局限与可能的发展趋势。
